\label{chap:83-15}

\section*{Mezcla CKM y violación de carga--paridad}

El sexteto de incidencia
\((b,p,q,h_6,o_8,t)=(4,5,7,6,8,3)\) fija la transformación racional
\begin{equation}\label{p12-83-c15-s1:eq:constantes-ckm-matriz-racional}
 M_{\rm CKM}=
 \begin{pmatrix}
 2&-5&28\\
 0&7&-25/2\\
 1&-20&-2/3
 \end{pmatrix},
 \qquad
 \det M_{\rm CKM}=-\frac{3857}{6}\ne0.
\end{equation}
Con
\[
 h_{\deg}=\frac{C^*_{\deg}}6,\qquad
 \Delta_{\deg}=\frac{180}{\pi}\Delta_4,
\]
se define
\[
 \begin{pmatrix}\theta_{12}\\\theta_{23}\\\theta_{13}\end{pmatrix}
 =M_{\rm CKM}
 \begin{pmatrix}A_{\deg}\\h_{\deg}\\\Delta_{\deg}\end{pmatrix},
 \qquad \delta_{{\rm CP},\deg}=9A_{\deg}.
\]
Para impedir toda ambigüedad de carta, se fija
\[
 s_{ij}:=\sin\!\left(\frac{\pi\theta_{ij}}{180}\right),\qquad
 c_{ij}:=\cos\!\left(\frac{\pi\theta_{ij}}{180}\right).
\]
La parametrización estándar produce una matriz unitaria \(V\), y el
invariante de Jarlskog es
\begin{equation}\label{p12-83-c15-s1:eq:constantes-jarlskog}
 J=c_{12}c_{23}c_{13}^2s_{12}s_{23}s_{13}
 \sin\!\left(\frac{\pi\delta_{{\rm CP},\deg}}{180}\right)
 =3.1189723326632974\times10^{-5}\ne0.
\end{equation}
La no anulación expresa violación CP en el modelo construido. El contraste
con el Particle Data Group se realiza después; no ajusta
\eqref{p12-83-c15-s1:eq:constantes-ckm-matriz-racional}. El sector PMNS requiere su propio
operador de mezcla y no se identifica con CKM ni con una fase diagonal de
rango uno.


\section{Construcción adimensional del escalar electrónico desde la incompatibilidad APP y el estado dodecafásico}
\label{p12-83-c15-s2:chap:biografia-electron-rev6}
\index{electrón!construcción holográfica}
\index{positrón!banda conjugada}
\index{masa!ley electrónica}

El electrón aparece en HMT como el punto en el que se ensamblan seis
construcciones anteriores y mutuamente tipadas: la incompatibilidad entre
las dos hojas aritméticas de APP, la orientación inducida por la involución
de la bisagra areal--volumétrica, dos cocientes enteros del censo APP--TPK,
el cierre
dodecafásico de \(\alpha_{\mathrm{HMT}}\), el defecto global
\(\Delta_4\) y el centro único del atlas nonádico. La banda de ruta añade
después la conjugación electrón--positrón. La carta en MeV y la comparación
con CODATA son posteriores a todas ellas.

Esta precedencia forma parte del contenido matemático: los factores, sus
tipos y su orden causal se determinan antes de abrir la carta tabulada, que
interviene únicamente en el contraste metrológico posterior.

\subsection{Escalar electrónico \(\mathfrak e_{\mathrm{HMT}}\) y grafo causal de sus factores APP–TPK, \(\alpha_{\mathrm{HMT}}\), \(\Delta_4\) y atlas nonádico}

La salida interna es el escalar adimensional
\begin{equation}
 \label{p12-83-c15-s2:eq:electron-holografico-rev6}
 \boxed{
 \mathfrak e_{\mathrm{HMT}}
 =
 \frac{\sqrt3}{4}
 \left[
 \exp\!\left(\frac{\varphi}{\pi^2}\right)
 -\frac{396}{18}\,\alpha_{\mathrm{HMT}}^3
 \right]
 \left[
 1+\frac{270}{18}\,\Delta_4
 \right].}
\end{equation}
Las identidades
\[
 \frac{396}{18}=22=27-5,
 \qquad
 \frac{270}{18}=15=\binom62=3\cdot5
 \]
permiten recuperar la escritura abreviada
\[
 \mathfrak e_{\mathrm{HMT}}
 =
 \frac{\sqrt3}{4}
 \left(e^{\varphi/\pi^2}-22\alpha_{\mathrm{HMT}}^3\right)
 (1+15\Delta_4).
\]
Aquí la letra \(e\) dentro de la exponencial es el número de Euler; el
subíndice de \(\mathfrak e_{\mathrm{HMT}}\) designa la sección electrónica.

El grafo causal de la ecuación es
\[
\begin{array}{ccl}
\text{proyectores APP en dimensión tres}
&\longrightarrow& \sqrt3/4,\\[2mm]
\text{calendario TPK y hojas areal--volumétrica}
&\longrightarrow&
F_{\mathrm{av}},\ \pi/\varphi^2,\
\varphi/\pi^2,\\[2mm]
\text{cierre dodecafásico}
&\longrightarrow&\alpha_{\mathrm{HMT}},\\[2mm]
\text{censo APP--TPK y familia electrónica}
&\longrightarrow&396,\ 18,\ 270,\ 22,\ 15,\ 5,\\[2mm]
\text{lecturas globales de }\pi,e\text{ y registro}
&\longrightarrow&\Delta_4,\\[2mm]
\text{atlas nonádico y anti-sección}
&\longrightarrow&(5,5),\
\text{banda electrón--positrón}.
\end{array}
\]
Ninguna flecha de este diagrama recibe como argumento la masa electrónica
de referencia.

\subsection{1.er factor: incompatibilidad exacta de las hojas APP}

Para \(d\geq2\), sea
\[
 \mathscr A_d=\{1,\ldots,9\}^d
\]
con medida uniforme, y sean
\[
 \Sigma_d(x)=\operatorname{dr}(x_1+\cdots+x_d),
 \qquad
 \Pi_d(x)=\operatorname{dr}(x_1\cdots x_d)
\]
los observables aditivo y multiplicativo. Sus esperanzas condicionales
ortogonales sobre \(\ell^2(\mathscr A_d)\) son
\(P_{\Sigma,d}\) y \(P_{\Pi,d}\). La
\cref{prop:app-no-conmutatividad} demuestra por enumeración exacta que
\[
 \left\|[P_{\Sigma,3},P_{\Pi,3}]\right\|
 =\frac{\sqrt3}{4}.
\]

\begin{proposition}[Contenido APP del prefactor electrónico]
\label{p12-83-c15-s2:prop:prefactor-app-electron-rev6}
El prefactor \(\sqrt3/4\) de
\cref{p12-83-c15-s2:eq:electron-holografico-rev6} es la norma exacta de la
no conmutatividad entre las hojas suma y producto en dimensión tres.
No es un factor geométrico introducido desde la fórmula electrónica.
\end{proposition}

\begin{proof}
La tabla conjunta de \(\Sigma_3\) y \(\Pi_3\) determina los ángulos
principales entre
\(\operatorname{Ran}P_{\Sigma,3}\) y
\(\operatorname{Ran}P_{\Pi,3}\). Si \(s\) es uno de sus valores
singulares, la contribución correspondiente a la norma del conmutador es
\(s\sqrt{1-s^2}\). La enumeración de las \(9^3\) palabras da como máximo
\[
 s^2(1-s^2)=\frac3{16},
 \qquad s^2=\frac14.
\]
Por tanto la norma máxima es \(\sqrt{3/16}=\sqrt3/4\), antes de
introducir \(\alpha_{\mathrm{HMT}}\), \(\Delta_4\), una unidad de energía
o una masa de referencia.
\end{proof}

La función física del factor queda así fijada: mide la incompatibilidad
máxima de las dos lecturas aritméticas sobre el mismo soporte ternario. El
electrón no hereda dos números independientes, uno aditivo y otro
multiplicativo; hereda la obstrucción exacta a hacer simultáneamente
compatibles las dos hojas. El prefactor conserva en la ley de masa la
memoria local del desacuerdo APP del que nace la ruta.

\subsection{2.º factor: la involución de intercambio entre \(\pi\) y \(\varphi\)}

El calendario primitivo de modos
\((\Sigma,\Pi,\Pi)\) induce, en las hojas areal y volumétrica, la matriz
\[
 F_{\mathrm{av}}
 =
 \begin{pmatrix}
 0&1\\
 1&1
 \end{pmatrix}.
\]
El \cref{p12-83-c16-s2:thm:descenso-necesario-fav} la obtiene contando las transiciones
\(\Sigma\to\Pi\), \(\Pi\to\Pi\) y \(\Pi\to\Sigma\), cada una una vez.
No se postula la matriz a partir de la razón áurea: sus autovalores
\(\varphi\) y \(-\varphi^{-1}\) aparecen después de haber construido la
incidencia.

En la memoria de segundo orden,
\(\operatorname{Sym}^2(F_{\mathrm{av}})\) tiene espectro
\[
 \varphi^2,\qquad -1,\qquad\varphi^{-2}.
\]
Así, \(\varphi^{-2}\) es el modo estable positivo de la bisagra. La marca
de clausura \(\pi\), aplicada una sola vez al retorno areal, produce
\[
 \frac{\pi}{\varphi^2}.
\]
Esta razón conserva el orden cierre--autoescala. La ley electrónica usa su
orientación contraria.

\begin{theorem}[Involución exacta de la bisagra electrónica]
\label{p12-83-c15-s2:thm:espejo-electron-rev6}
Sean
\[
 \mathscr R(x,y)=\frac{x}{y^2},
 \qquad
 \mathsf J(x,y)=(y,x).
\]
Entonces \(\mathsf J^2=I\) y
\[
 \mathscr R(\pi,\varphi)=\frac{\pi}{\varphi^2},
 \qquad
 (\mathscr R\circ\mathsf J)(\pi,\varphi)
 =\frac{\varphi}{\pi^2}.
\]
Por consiguiente,
\[
 \frac{\varphi}{\pi^2}
 =
 \frac{1}{
 \varphi^3\bigl(\pi/\varphi^2\bigr)^2},
 \qquad
 e^{\varphi/\pi^2}
 =
 \exp\!\left[
 \frac{1}{
 \varphi^3\bigl(\pi/\varphi^2\bigr)^2}
 \right].
\]
\end{theorem}

\begin{proof}
La primera pareja de identidades se obtiene aplicando
\(\mathscr R\) antes y después de la involución \(\mathsf J\). Para la
reparametrización,
\[
 \varphi^3
 \left(\frac{\pi}{\varphi^2}\right)^2
 =\frac{\pi^2}{\varphi},
\]
y su inversa es \(\varphi/\pi^2\). La última igualdad se sigue al aplicar
la exponencial.
\end{proof}

La presencia simultánea de \(\pi\) y \(\varphi\) tiene, por tanto, un
significado estructural preciso. \(\pi/\varphi^2\) lee el cierre areal
sobre la memoria volumétrica cuadrática; \(\varphi/\pi^2\) lee el
autocrecimiento interior sobre el cierre cuadrático. No son dos cocientes
parecidos elegidos por conveniencia. Son las dos orientaciones de un mismo
lector, intercambiadas por una involución exacta. El término
\(e^{\varphi/\pi^2}\) publica en la sección electrónica la orientación
interior de esa bisagra.

\subsection{3.er y 4.º factores: los enteros 22, 15 y el centro 5}

El censo de la carta APP empleada por la familia electrónica proporciona
tres sumas enteras:
\[
 \Sigma(\mathrm{APP})=396,
 \qquad
 \Sigma(\mathrm{centro}\ 2\times2)=18,
 \qquad
 \Sigma(\mathrm{canal})=270.
\]
La primera admite la descomposición por anillos
\[
 396=18+72+108+198.
\]
El núcleo \(2\times2\) es, por tanto, una subestructura explícita de la
carta completa, y el entero 270 pertenece al canal global del mismo censo
APP--TPK.
La suma de los anillos pares es
\[
 18+108=126,
 \qquad
 \frac{396}{126}=\frac{22}{7}.
\]
Este último cociente es el sello racional de cierre de la carta, no una
identificación de \(22/7\) con el número \(\pi\) generado por el lector
coinductivo. Sí muestra que el numerador \(22\) ya ocupa una función
interna antes de aparecer en la ley electrónica.

\begin{theorem}[Procedencia interna de los coeficientes electrónicos]
\label{p12-83-c15-s2:thm:coeficientes-electron-rev6}
En la familia electrónica fijada por la carta APP y su canal global,
\[
 \boxed{
 22=\frac{\Sigma(\mathrm{APP})}
 {\Sigma(\mathrm{centro})}
 =\frac{396}{18}=27-5,}
\]
\[
 \boxed{
 15=\frac{\Sigma(\mathrm{canal})}
 {\Sigma(\mathrm{centro})}
 =\frac{270}{18}=3\cdot5=\binom62.}
\]
En particular, los coeficientes se determinan sin consultar
\(Z_0\), la masa del electrón ni una carta de unidades.
\end{theorem}

\begin{proof}
Las dos primeras igualdades son divisiones exactas del censo finito:
\[
 396=22\cdot18,
 \qquad
 270=15\cdot18.
\]
Al definir
\[
 p:=27-22,
\]
se obtiene \(p=5\), y entonces
\[
 15=3p=3\cdot5=\binom62.
\]
Todas las cantidades pertenecen al dominio entero de la carta y de la
familia antes de la transducción metrológica.
\end{proof}

\subsubsection{Selector de persistencia en el retorno central}

La misma pareja posee una construcción dinámica que fija además su signo sin
consultar una masa. Las cuatro orientaciones de la semilla central
\((5,5)\) regresan por primera vez a esa celda tras veintisiete pasos. El
espacio tipado del retorno es
\[
 \mathcal K_e=\mathbb Z/9\mathbb Z\times\mathbb F_3,
 \qquad |\mathcal K_e|=27.
\]
Por otra parte, la microfibra de \(\pi\) contiene las cinco regiones
genealógicamente distintas
\[
 \mathscr R_\pi=3_{++}+1_{--}+1_{\perp}.
\]
El calibre regional \(\operatorname{diag}_c\) y, dentro de una fase repetida,
el orden interno de \(U_6\), definen la inyección
\[
 \iota:\mathscr R_\pi\hookrightarrow\mathcal K_e
\]
mediante los cinco lugares distintos
\[
 (4,0),\quad(5,0),\quad(6,0),\quad(6,1),\quad(6,2).
\]
Los dos primeros corresponden, respectivamente, a la región transversal y
al retorno mutado; los tres últimos separan las regiones positivas por su
hoja trítica. La palabra visible común no interviene como inversa del mapa.

\begin{theorem}[Selector interno de los coeficientes electrónicos]
\label{p12-83-c15-s2:thm:selector-persistencia-electron-rev7}
Fijados el retorno central y las cinco regiones tipadas, los coeficientes de
la ecuación electrónica son
\[
 \boxed{
 n_e=-\bigl|\mathcal K_e\setminus\iota(\mathscr R_\pi)\bigr|=-22,
 \qquad
 m_e^{\rm tor}=\bigl|\mathscr R_\pi\times\mathbb F_3\bigr|=+15.}
\]
El primer signo registra el déficit de frontera; el segundo, la memoria
torsional retenida.
\end{theorem}

\begin{proof}
Los cinco lugares anteriores son distintos, de modo que \(\iota\) es
inyectiva y
\[
 |\mathcal K_e\setminus\iota(\mathscr R_\pi)|=27-5=22.
\]
Cada una de las cinco regiones conserva tres hojas tríticas, luego
\[
 |\mathscr R_\pi\times\mathbb F_3|=5\cdot3=15.
\]
La regla de orientación asigna signo negativo a los estados retirados de la
cáscara y positivo a la memoria retenida. La construcción sólo abre el
centro, el retorno de veintisiete estados y las identidades regionales; no
abre una masa observada, una impedancia ni una firma calibrada.
\end{proof}

El número \(22\) mide cuántos núcleos centrales de peso \(18\) recompone la
suma APP \(396\). El número \(15\) mide cuántos de esos mismos núcleos
recompone el canal \(270\), y simultáneamente cuenta los pares de una
hexada. El número \(5\) no es un parámetro adicional:
\[
 5=27-22,\qquad 15=3\cdot5.
\]
Las tres escrituras muestran el mismo cierre desde la profundidad 27,
desde la normalización central y desde la combinatoria hexádica.

\subsubsection{Derivación interna de los coeficientes 22 y 15 mediante los censos APP--TPK}

La exploración histórica localizó primero \(p=5\) mediante una comparación
finita con \(Z_0\), y de allí escribió
\((27-p,3p)=(22,15)\). Ese procedimiento conserva valor como historia de
descubrimiento y como control externo. No es, sin embargo, la demostración
que gobierna la ley: el recuento posterior
\[
 \frac{396}{18}=22,
 \qquad
 \frac{270}{18}=15
\]
construye ambos enteros dentro de APP--TPK y recupera \(p=5\) como
consecuencia. La historia explica cómo se encontró la pareja; el censo
explica por qué pertenece a la ecuación.

\subsection{5.º factor: el cubo del cierre \(\alpha_{\mathrm{HMT}}\)}

El cierre de \(\alpha_{\mathrm{HMT}}\) precede a Mecánica Dimensional.
Como se establece en el \cref{chap:alpha}, los doce bloques en base mil
de \(\pi\), \(e\), \(\varphi\) y del sello \(K\) se combinan mediante el
acarreo reversible
\[
 P_m+E_m-\Phi_m-K_m+c_{m+1}
 =A_m+1000c_m,
 \qquad c_{13}=c_1=0.
\]
La salida es
\[
 A_\alpha=
 (007,297,352,569,283,800,997,285,105,472,380,663)
\]
y, por concatenación, la ventana dodecafásica exacta
\[
 \alpha_{12}=\pi_{12}(\alpha_{\mathrm{HMT}})
 =
 \frac{
 7297352569283800997285105472380663
 }{10^{36}}.
\]
El telescopado de los acarreos demuestra la identidad entera completa; la
prolongación compatible construye
\(\alpha_{\mathrm{HMT}}=\varprojlim_N A^{(N)}\), y \(\alpha_{12}\) es su
proyección racional de treinta y seis cifras. La
comparación con la constante de estructura fina se realiza después y no
interviene en su construcción.

La ley electrónica toma la tercera potencia
\[
 \alpha_{\mathrm{HMT}}^3.
\]
Matemáticamente, se trata de la evaluación multiplicativa de la tercera
potencia simétrica del mismo escalar ya cerrado; no se introducen tres
copias metrológicas ni tres parámetros diferentes. Su contribución es
\[
 22\alpha_{\mathrm{HMT}}^3
 =
 0.00000854906599200551727014979807598228\ldots.
\]

\begin{proposition}[Estatuto del canal cúbico electrónico]
\label{p12-83-c15-s2:prop:alpha-cubica-electron-rev6}
El término \(22\alpha_{\mathrm{HMT}}^3\) de la ley electrónica depende
únicamente del cierre previo de \(\alpha_{\mathrm{HMT}}\), de la
multiplicación interna y del entero \(22\) del
\cref{p12-83-c15-s2:thm:coeficientes-electron-rev6}. No depende de una lectura de la
masa electrónica.
\end{proposition}

\begin{proof}
\(\alpha_{12}\) es el racional fijado por la identidad de concatenación y
acarreo, y la familia compatible de ventanas determina la sección
\(\alpha_{\mathrm{HMT}}\). La proyección \(\alpha_{12}^3\) certifica las
cifras impresas de \(\alpha_{\mathrm{HMT}}^3\).
El coeficiente \(22\) está determinado por \(396/18\). La composición de
ambas operaciones contiene sólo datos construidos con anterioridad a la
carta de energía.
\end{proof}

Esta proposición fija exactamente lo que se usa en el electrón. La potencia
tres es el grado cúbico de este canal en la ley publicada. No es necesario
postular aquí una regla universal según la cual toda aparición
\(\alpha^d\) represente automáticamente una dimensión \(d\); una
afirmación universal de ese tipo exigiría su propio teorema de
multiplicatividad entre grados. La genealogía electrónica queda completa
sin convertir esa generalización en premisa.

\subsection{6.º factor: torsión mínima y defecto global \(\Delta_4\)}
\index{torsión mínima!\(\Delta_4\)}

Las denominaciones \emph{torsión mínima} y \emph{defecto global}
designan aquí un único invariante, no dos correcciones distintas. El objeto
empleado por la ley de masas es
\begin{equation}
 \label{p12-83-c15-s2:eq:delta4-electron-rev6}
 \boxed{
 \Delta_4
 =
 \frac{\pi-e\log\pi}{270}
 \left(1+\frac{\pi}{729}\right).}
\end{equation}
Sus entradas son las lecturas HMT ya fijadas de \(\pi\) y \(e\), el canal
global \(270\) y la celda interior \(729\). Por consiguiente,
\[
 \Delta_4
 =
 0.00011119642269214005405113478612067909\ldots.
\]
La primera razón mide el defecto firmado
\(\pi-e\log\pi\) en unidades del canal 270; el segundo factor acopla esa
razón a la celda interior mediante \(\pi/729\). Esta lectura no añade
parámetros a la fórmula: separa sus dos operaciones constitutivas.
El \cref{p3-20260826:chap:medida-accion-positiva} muestra que es una entrada matemática
pre-dimensional y la distingue del operador que eventualmente la realiza
como cambio de hoja.

\begin{proposition}[Globalidad del defecto electrónico]
\label{p12-83-c15-s2:prop:delta4-global-electron-rev6}
\(\Delta_4\) es un invariante global anterior a la especie. En la ley
electrónica entra mediante
\[
 1+15\Delta_4
 =
 1.00166794634038210081076702179181019\ldots.
\]
No es un corrector elegido después de conocer el residuo de la masa del
electrón.
\end{proposition}

\begin{proof}
La \cref{p12-83-c15-s2:eq:delta4-electron-rev6} no contiene una etiqueta de especie, una
masa ni una unidad. El entero \(15\) ya está construido por \(270/18\).
Por tanto la composición \(1+15\Delta_4\) queda determinada por los datos
globales y el censo interno.
\end{proof}

Debe distinguirse \(\Delta_4\) de los otros bloques del carácter de masa y de
las otras nociones de torsión. En particular, no es \(s_{270}\), ni el
invariante cuadrático \(135\Delta_4^2\), ni el residuo angular
\(\varepsilon_{\rm res}\), ni la torsión geométrica de Cartan--Holst.
Sustituir uno por otro repetiría o cambiaría el nivel de la corrección. En el
electrón, \(\Delta_4\) conserva la memoria global del refinamiento; el factor
\(15\) declara cuántas veces la familia electrónica la lee. El calificativo
\emph{mínima} conserva el nombre histórico de este invariante y no se usa
como sustituto de un teorema variacional de minimalidad.

\subsection{Unicidad de la celda central (5,5) como soporte genealógico del electrón}

Sea
\[
 \mathcal C_{81}=\{1,\ldots,9\}^2
\]
el atlas nonádico y
\[
 \rho(r,c)=(10-r,10-c)
\]
su inversión central. El \cref{thm:censo-atlas-81} prueba que el atlas
contiene una única celda de generación cero y que ésta porta la etiqueta
posicional \(e_{\mathrm{core}}\).

\begin{proposition}[Unicidad del centro electrónico]
\label{p12-83-c15-s2:prop:centro-electron-rev6}
La involución \(\rho\) posee un único punto fijo:
\[
 \boxed{(r,c)=(5,5).}
\]
Su unicidad es combinatoria y no depende del valor de
\(\mathfrak e_{\mathrm{HMT}}\).
\end{proposition}

\begin{proof}
La ecuación \((r,c)=\rho(r,c)\) equivale a
\[
 2r=10,\qquad2c=10.
\]
En \(\{1,\ldots,9\}^2\) tiene la única solución \(r=c=5\). La evaluación
de la masa no aparece en el argumento.
\end{proof}

\begin{corollary}[Universalidad de la sección central]
\label{p12-83-c15-s2:cor:universalidad-electron-central-rev6}
Todo automorfismo admisible del atlas que entrelace la inversión central
preserva la celda \((5,5)\).
\end{corollary}

\begin{proof}
Si \(f\rho=\rho f\) y \(\rho(5,5)=(5,5)\), entonces
\(\rho f(5,5)=f\rho(5,5)=f(5,5)\). Por tanto \(f(5,5)\) es otro punto fijo
de \(\rho\); la unicidad obliga a \(f(5,5)=(5,5)\).
\end{proof}

Ésta es la razón estructural por la que la clase electrónica es la misma en
todas las realizaciones que preservan el atlas y su inversión: no se repite
por casualidad un decimal local, sino que toda carta admisible debe
transportar el mismo punto fijo y la misma cadena de invariantes. La lectura
física de esta unicidad identifica cada electrón con una realización de esa
clase central; la proposición demuestra la universalidad interna del atlas.

En la familia electrónica fijada por la construcción, este punto es el anclaje
geométrico de la sección. Conviene no confundir cuatro objetos:
\[
 (5,5),\qquad
 e_{\mathrm{core}},\qquad
 v_e=(0,0,0,0,0),\qquad
 \text{la firma cruda de la celda central}.
\]
El primero es una posición; el segundo, su etiqueta interna; el tercero,
el origen elegido en el retículo fino de acción; el cuarto conserva los
datos crudos del atlas y no tiene por qué ser el vector nulo. Su
coordinación explica por qué el electrón ocupa el centro sin borrar la
genealogía de las representaciones que lo describen.

\paragraph{Rutas electrónicas orientadas \(e_{++},e_{--}\), envolvente \(e_{\mathrm{env}}\), origen de acción y firmas de la celda central}
La celda central porta dos historias orientadas de doce eventos. Escribiendo
\(v^2\) para la concatenación de dos copias de una sextupla, la ruta activa
de la orientación \({++}\), su historia opuesta y la envolvente histórica
son, respectivamente,
\[
\begin{aligned}
 e_{\rm act}=e_{++}
   &=(2,2,5,-4,-3,-2)^2,\\
 e_{--}
   &=(4,3,2,-2,-2,-5)^2,\\
 e_{\rm env}
   &=(4,3,5,-4,-3,-5)^2.
\end{aligned}
\tag{E-centro-1}\label{p12-83-c15-s2:eq:electron-rutas-finas-rev8}
\]
Las dos historias publican la misma palabra
\[
 \tau_{12}=\texttt{+++---+++---},
\]
pero no son la misma ruta. La envolvente toma el máximo absoluto
coordenada a coordenada conservando el signo; registra amplitudes, pero
aplana el orden genealógico y no sustituye a \(e_{\rm act}\).

Los otros lectores del centro viven en codominios distintos:
\[
\begin{array}{rcll}
 v_e&=&(0,0,0,0,0),
   &\text{origen afín del retículo de acción},\\
 R(5,5)&=&(24,0,12,0,-12),
   &\text{firma cruda APP de la celda},\\
 \operatorname{Sig}^{\Gamma_{108}}_{\rm masa}
   &=&(-12,-4,12,-6,12),
   &\text{proyección par de masa},\\
 \operatorname{Sig}^{\Gamma_{108}}_{\rm carga}
   &=&(0,0,0,0),
   &\text{proyección impar de carga}.
\end{array}
\tag{E-centro-2}\label{p12-83-c15-s2:eq:electron-firmas-tipadas-rev8}
\]
Que la palabra de carga tenga firma impar nula no convierte la ruta en
cero, ni identifica la firma cruda con el origen afín. Posición, ruta
activa, ruta opuesta, envolvente, origen y ambas firmas son objetos
coordinados, no seis notaciones para un mismo vector.

\subsection{Órbita conjugada electrón–positrón sobre la banda de rutas orientadas}

La identidad electrónica no termina en el punto central ni en el valor de
masa. La conjugación material se realiza sobre palabras dodecafásicas. Sea
\[
 C_M(\mathbf t)=-\operatorname{rev}(\mathbf t)
\]
la anti-sección orientada y sea
\(\chi\in\{-1,+1\}\) la orientación de carga. La banda es la clase
\[
 [(\mathbf t,\chi)]_M
 =
 \{(\mathbf t,\chi),(C_M\mathbf t,-\chi)\}
\]
de la \cref{def:banda-dodecafasica-ruta}. La banda conserva sus componentes
pares y presenta dos secciones con componentes impares opuestas.

\begin{theorem}[Electrón y positrón como secciones de una banda]
\label{p12-83-c15-s2:thm:electron-positron-biografia-rev6}
Para el carácter de banda del
\cref{thm:principio-conjugacion-masa-banda},
\[
 m(C_M\mathbf t,-\chi)=m(\mathbf t,\chi).
\]
En el canal electrónico,
\[
 (n,\chi)=(-22,+1)
 \quad\longleftrightarrow\quad
 (+22,-1)
\]
y ambos pares producen
\[
 \chi n=-22.
\]
Por tanto electrón y positrón tienen la misma masa y orientaciones de
carga opuestas.
\end{theorem}

\begin{proof}
Si \(E_+\) y \(E_-\) son, respectivamente, las componentes par e impar del
carácter respecto de \(C_M\), entonces
\[
 E_+(C_M\mathbf t)=E_+(\mathbf t),
 \qquad
 E_-(C_M\mathbf t)=-E_-(\mathbf t).
\]
De aquí,
\[
 E_{\mathrm{band}}(C_M\mathbf t,-\chi)
 =
 E_{\mathrm{band}}(\mathbf t,\chi),
\]
y la igualdad de masa se obtiene al aplicar la exponencial. En el término
electrónico,
\[
 (+1)(-22)=(-1)(+22)=-22.
\]
\end{proof}

El cierre finito de la \cref{prop:cierre-antiseccion-ct108} refuerza esta
interpretación: sobre las ochenta y una rutas CT--108, ni el signo puro ni
la reversión pura preservan la taxonomía, mientras
\(-\operatorname{rev}\) sí la preserva, con
\[
 81=54+9+9+9.
\]
La antipartícula no es, por tanto, una masa negativa ni una copia
independiente. Es la otra sección orientada de la misma banda: conserva los
invariantes pares que determinan la masa e invierte los impares que
determinan carga y orientación.

\subsection{Etapa basal del cierre electrónico}

\begin{theorem}[Composición holográfica de la sección electrónica]
\label{p12-83-c15-s2:thm:cierre-electron-holografico-rev6}
Fijados los objetos construidos en las secciones anteriores, la sección
electrónica adimensional es
\[
 \mathfrak e_{\mathrm{HMT}}
 =
 \left\|[P_{\Sigma,3},P_{\Pi,3}]\right\|
 \left[
 \exp\bigl((\mathscr R\circ\mathsf J)(\pi,\varphi)\bigr)
 -\frac{\Sigma(\mathrm{APP})}
 {\Sigma(\mathrm{centro})}\,
 \alpha_{\mathrm{HMT}}^3
 \right]
 \left[
 1+
 \frac{\Sigma(\mathrm{canal})}
 {\Sigma(\mathrm{centro})}\,
 \Delta_4
 \right].
\]
Equivalentemente,
\[
 \boxed{
 \mathfrak e_{\mathrm{HMT}}
 =
 \frac{\sqrt3}{4}
 \left(e^{\varphi/\pi^2}
 -22\alpha_{\mathrm{HMT}}^3\right)
 (1+15\Delta_4).}
\]
Su evaluación es
\[
 \boxed{
 \mathfrak e_{\mathrm{HMT}}
 =
 0.51099892248806607250987087719134943763\ldots.}
\]
Este escalar es la etapa basal \(\mathfrak e_e^{(0)}\), no la salida
electrónica rectora. La clausura bidireccional beta publica
\(0.510998950690000808608\ldots\,\mathrm{MeV}\).
\end{theorem}

\begin{proof}
La \cref{p12-83-c15-s2:prop:prefactor-app-electron-rev6} aporta
\[
 \left\|[P_{\Sigma,3},P_{\Pi,3}]\right\|=\sqrt3/4.
\]
El \cref{p12-83-c15-s2:thm:espejo-electron-rev6} aporta
\[
 (\mathscr R\circ\mathsf J)(\pi,\varphi)=\varphi/\pi^2.
\]
El \cref{p12-83-c15-s2:thm:coeficientes-electron-rev6} aporta
\[
 \frac{\Sigma(\mathrm{APP})}{\Sigma(\mathrm{centro})}=22,
 \qquad
 \frac{\Sigma(\mathrm{canal})}{\Sigma(\mathrm{centro})}=15.
\]
El cierre dodecafásico aporta la ventana racional \(\alpha_{12}\) de la
sección coinductiva \(\alpha_{\mathrm{HMT}}\), y
\cref{p12-83-c15-s2:eq:delta4-electron-rev6} aporta el defecto global. La sustitución
produce la fórmula abreviada. La evaluación sucesiva de
\[
\begin{aligned}
 e^{\varphi/\pi^2}
 &=1.17814494259630678081160095680888910\ldots,\\
 22\alpha_{\mathrm{HMT}}^3
 &=0.00000854906599200551727014979807598228\ldots,\\
 1+15\Delta_4
 &=1.00166794634038210081076702179181019\ldots
\end{aligned}
\]
da el valor declarado.
\end{proof}

El teorema es holográfico en un sentido verificable: cada factor local
retiene una capa distinta del sistema completo. El conmutador retiene la
dualidad de hojas; la exponencial retiene la orientación interior de la
bisagra \(\pi\)--\(\varphi\); \(22\) y \(15\) retienen el censo global y
su núcleo central; \(\alpha_{\mathrm{HMT}}^3\) retiene el cierre fino; y
\(\Delta_4\) retiene el refinamiento global. El centro determina dónde se
ancla la sección y la banda determina cómo sobrevive bajo conjugación.
El valor no reemplaza esta estructura: es su evaluación común.

\subsubsection{Desdoblamiento de acción anterior y posterior al retorno}
\label{p12-83-c15-s2:sec:electron-desdoblamiento-two-way-rev9}

La realización regional de la coordenada angular conjugada \(C^*\) distingue dos coordenadas sobre la
misma recta interna de acción. La coordenada anterior al retorno es \(H_5\) y
la posterior es
\[
\etaRet
 =H_5-\frac{90}{\pi}\mathscr C_\pi(\alpha_{\rm HMT}).
\]
La notación \(\etaRet\) es la coordenada sexagesimal canónica definida en la
rama de acción por \(\etaRet=\Cstar/2-\alphaAngle\). La igualdad anterior se
evalúa en esa carta; no introduce una segunda coordenada de retorno.
Por tanto, las cartas aditiva y multiplicativa del desdoblamiento orientado
son
\begin{equation}
 \boxed{
 \delta_{2w}:=H_5-\etaRet
 =\frac{90}{\pi}\mathscr C_\pi(\alpha_{\rm HMT}),
 \qquad
 R_{\rm act}:=\frac{\hbar^{\rm pre}_{\rm HMT}}
 {\hbar^{\rm ret}_{\rm HMT}}
 =\frac{H_5}{\etaRet}.}
 \label{p12-83-c15-s2:eq:electron-desdoblamiento-two-way-rev9}
\end{equation}
Ambas coordenadas proceden de un único transporte pre-retorno/post-retorno.
No son dos constantes de Planck: el factor común
\(10^{-34}\mathcal U_S\) fija la recta dimensional y cancela exactamente en
\(R_{\rm act}\). La elipse de acción realiza después
\(\operatorname{Area}(\theta)=\hbar_{\rm HMT}^{\rm ret}\theta\), de modo que
el desdoblamiento antecede a la lectura de holonomía.

El acoplamiento
\begin{equation}
 \kappa_e^{\beta}
 =80-54\alpha_{\rm HMT}+6\Delta_4,
 \qquad
 \mathfrak e_e^{\beta}
 =\mathfrak e_{\rm HMT}R_{\rm act}^{\kappa_e^{\beta}}
 =0.5109989506900008086\ldots
 \label{p12-83-c15-s2:eq:electron-beta-two-way-rev9}
\end{equation}
es la clausura cuantitativa bidireccional del electrón. Los dos lectores
independientes del registro electrónico construidos en el capítulo siguiente
seleccionan, sin consultar la masa de contraste,
\[
 K_D(\Gamma_e)=K_\Omega(\Gamma_e)=(80,54,6),
\]
y derivan exactamente \(\kappa_e^\beta\). El entero \(80\) procede de la
primera coordenada de esos lectores; la palabra \(w_e=201101\) conserva su
residencia propia en el canal de la constante de Euler.

\begin{statusbox}[title={Clausura interna del refinamiento bidireccional}]
La identidad \eqref{p12-83-c15-s2:eq:electron-desdoblamiento-two-way-rev9}, la reducción
\((80,54,6)\) y la ecuación
\eqref{p12-83-c15-s2:eq:electron-beta-two-way-rev9} constituyen resultados internos exactos
y target--free sobre el registro electrónico. La propagación a fibras de
rango mayor se formula mediante los operadores de multisección.
\end{statusbox}
